\documentclass[12pt]{article}
\usepackage[T1]{fontenc}
\usepackage[utf8]{inputenc}
\usepackage{lmodern}
\usepackage{textcomp}
\usepackage{enumitem}
\usepackage{geometry}
\geometry{margin=1in}
\begin{document}
\begin{center}
Answer Key - Physics - Undergraduate Level Assessment 11 \
\small (Answers are ordered exactly as in the question paper.)
\end{center}

\section*{Multiple Choice}
\begin{enumerate}[label=\arabic*.]
\item \textbf{Answer:} B
\\ \textit{Options:} A) Magnetic susceptibility; B) Hall coefficient; C) Electrical resistivity; D) Thermal conductivity
\\ \textit{Note:} The Hall coefficient indicates the nature of charge carriers in a doped semiconductor by determining whether they are positive (holes) or negative (electrons) based on the direction of the induced voltage in a magnetic field.
\item \textbf{Answer:} A
\\ \textit{Options:} A) 5; B) 4; C) 3; D) 2
\\ \textit{Note:} The quantum number m\_l can take on values ranging from -l to +l, inclusive, so for l = 2, the possible values of m\_l are -2, -1, 0, +1, and +2, resulting in a total of 5 allowed values.
\item \textbf{Answer:} B
\\ \textit{Options:} A) 0.50 c; B) 0.60 c; C) 0.70 c; D) 0.80 c
\\ \textit{Note:} The correct answer is 0.60 c because, according to the principles of special relativity, length contraction occurs for an object moving relative to an observer, and the relationship between the proper length (1.00 m) and the contracted length (0.80 m) can be calculated using the formula L = L$_{0}$√(1 - v²/c²), which yields v = 0.60 c.
\item \textbf{Answer:} A
\\ \textit{Options:} A) Planck's constant; B) Boltzmann's constant; C) The Rydberg constant; D) The speed of light
\\ \textit{Note:} De Broglie's hypothesis states that a particle's wavelength is inversely proportional to its momentum, with Planck's constant serving as the proportionality factor, establishing the fundamental relationship between wave and particle properties in quantum mechanics.
\item \textbf{Answer:} A
\\ \textit{Options:} A) 4; B) 5; C) 6; D) 20
\\ \textit{Note:} The angular magnification of a refracting telescope is given by the formula M = -f\_o / f\_e, where f\_o is the focal length of the objective lens and f\_e is the focal length of the eyepiece; with the eyepiece lens having a focal length of 20 cm and the objective lens having a focal length of 80 cm (derived from the total separation of 100 cm), the magnification calculates to 4.
\end{enumerate}

\section*{True / False}
\begin{enumerate}[label=\arabic*.]
\setcounter{enumi}{5}
\item \textbf{Answer:} True
\\ \textit{Options:} A) True; B) False
\\ \textit{Note:} Fermions are particles that have half-integer spin and, according to the principles of quantum mechanics, their wave functions must be antisymmetric under the exchange of two identical fermions, which leads to the Pauli exclusion principle that states no two fermions can occupy the same quantum state simultaneously.
\item \textbf{Answer:} False
\\ \textit{Options:} A) True; B) False
\\ \textit{Note:} The statement is false because, according to special relativity, a particle traveling at 0.60 c will experience time dilation, meaning it will decay after traveling a shorter distance than 288 m in the lab frame due to its relativistic effects.
\item \textbf{Answer:} True
\\ \textit{Options:} A) True; B) False
\\ \textit{Note:} The angular magnification of a refracting telescope is calculated as the ratio of the focal lengths of the objective lens to the eyepiece lens, and if the focal length of the eyepiece is 20 cm, the objective lens must have a focal length of 80 cm to achieve an angular magnification of 4.
\item \textbf{Answer:} False
\\ \textit{Options:} A) True; B) False
\\ \textit{Note:} The mean kinetic energy of conduction electrons in metals is not necessarily high because, while they are free to move, their kinetic energy is primarily influenced by temperature and the effective mass of the electrons rather than the number of degrees of freedom.
\item \textbf{Answer:} True
\\ \textit{Options:} A) True; B) False
\\ \textit{Note:} The statement is true because, according to the Stefan-Boltzmann Law, the power radiated per unit area of a blackbody is proportional to the fourth power of its absolute temperature (P ∝ $T^4$), so increasing the temperature by a factor of 3 results in an increase in energy radiated by a factor of $3^4$, which equals 81.
\end{enumerate}

\section*{Long Form Response}
\begin{enumerate}[label=\arabic*.]
\setcounter{enumi}{10}
\item \textbf{Answer:} In the ground state of neutral nitrogen (Z = 7), the total spin quantum number of the electrons can be calculated by examining its electron configuration, which is 1 $s^2$ 2 $s^2$ 2 $p^3$. The two electrons in the 1 s and 2 s subshells are paired and contribute no net spin, while the three electrons in the 2 p subshell each have a spin of ½. According to Hund's rule, these three 2 p electrons will occupy separate orbitals with parallel spins, resulting in a total spin of ½ + ½ + ½ = 3/2. This total spin quantum number is significant as it influences the magnetic properties of the atom and plays a crucial role in determining the atom's behavior in chemical bonding and interactions.
\\ \textit{Note:} The answer correctly identifies the total spin quantum number of nitrogen's electrons by analyzing its electron configuration, noting that paired electrons contribute zero net spin while unpaired electrons in the 2 p subshell follow Hund's rule to maximize total spin, resulting in a total spin of 3/2.
\item \textbf{Answer:} An atom with filled n = 1 and n = 2 energy levels has an electron configuration of 1 $s^2$ 2 $s^2$ 2 p⁶. The n = 1 level can hold a maximum of 2 electrons in the 1 s subshell, while the n = 2 level can hold a maximum of 8 electrons, distributed among the 2 s and 2 p subshells (2 in 2 s and 6 in 2 p). Therefore, the total number of electrons in the atom is 2 (from n = 1) + 8 (from n = 2), resulting in 10 electrons overall. This filled configuration contributes to the atom's stability and can determine its chemical properties, as the filled outer shells typically indicate that the atom is a noble gas, which is chemically inert.
\\ \textit{Note:} correct because it accurately describes the maximum electron capacity of the filled n = 1 and n = 2 energy levels, leading to a total of 10 electrons, which reflects the atom's stable electron configuration.
\end{enumerate}

\vfill
\noindent\textit{End of Answer Key}
\end{document}